% ===== §10 =====
\section{Optimizing Practice Is Not Ethical Practice}
\label{p25:sec:ethics}

\noindent
This section is the pivot of the criticism and, in a sense, of the paper, for it states in its sharpest form the claim the title's question was aimed at from the start: that morality cannot be defined as an optimization problem, and that an optimizing practice, however well it succeeds on its own scale, is not thereby an ethical practice. The preceding sections argued that optimization is ill-posed in intimate relation and that its imposition does damage. This one argues something orthogonal to both, and independent of both: that even where an optimization is perfectly well-posed and perfectly successful, reaching a genuine maximum of a genuine objective, it may still fail as an ethics, because the criterion of the optimal and the criterion of the good are different criteria, answering different questions, and the optimal is not the good. The claim thus holds even in the shipping-fleet domain where the four conditions are met; it is not a claim about ill-posedness at all, but about the gap between maximizing and acting well, which no well-posedness closes.

\subsection*{The orthogonality of the optimal and the good}

The series has met a structure of this kind before, and naming the earlier instance sharpens the present one and shows it is not improvised for the occasion. In the paper on the sustainability of generative relation, poetry was shown to satisfy a formal, dynamical criterion, the criterion of the good cycle, of endless non-settling circulation; and it was shown in the same breath that satisfying this criterion does not make a symbolic circulation \emph{good}, since the aesthetics of fascism and the discourse of the cult satisfy the very same dynamical criterion and are evil \parencite{paperix}. The lesson drawn there was that the dynamical criterion and the justice criterion are orthogonal: the first asks whether a circulation can sustain itself, the second whether it ought to and at whose expense, and a circulation may pass the first and fail the second, so that goodness requires the justice criterion as an independent condition which the dynamics can never supply. The present claim has the same shape, with optimization in the place of dynamics. The criterion of the optimal answers one question: whether, on a given objective, a given option is the best available, whether it sits at the maximum. The criterion of the good answers another: whether the option ought to be chosen, whether it wrongs anyone, whether it honors what the agent is bound to honor. These two are orthogonal. An option may be optimal and wrong; an option may be right and suboptimal; and success on the first criterion is no success whatever on the second, since they do not measure one thing along a common scale but pose different questions that a single answer cannot jointly settle. To have found the maximum is to have answered the first question and not so much as touched the second.

\begin{casebox}{the optimal betrayal}
Suppose a man faces a choice well-posed enough to optimize: the relevant goods admit of proxies, the objective is stable, the scale is available, and he himself is settled in his preferences. He computes, correctly, that abandoning a friend who has become burdensome, at a moment when the friend most depends on him, maximizes his objective, and even, on a plausible specification, maximizes an impartial aggregate of well-being, since his own gain is large and the friend will in time recover. The optimization is impeccable; the maximum is genuine; on its own scale nothing better is available. And the act is a betrayal. The point is not that he has optimized the wrong objective, which would return us to the earlier impossibilities; grant that the objective is as good as an objective gets. The point is that optimality and rightness have come apart, that the very act which is best on the scale is the act that wrongs the one to whom he was bound, and that no refinement of the maximization repairs this, because the wrong is not a quantity the objective failed to include but a feature of the act that the whole dimension of maximizing cannot see. He did the optimal thing, and he ought not to have done it, and both are true at once.
\end{casebox}

\subsection*{The resistance of morality to the optimizing form}

That the two criteria are orthogonal is not a brute fact but follows from the kind of thing morality is, and three features of the moral show why it resists the optimizing form in particular, each naming a place where the case above bites. The first is the \emph{unconditionality} of certain moral requirements. An obligation of the kind Kant placed at the center of ethics does not present itself as one good among others to be weighed in a sum and outweighed should the sum come out otherwise; it presents itself as binding regardless of the balance, a constraint on what may be done and not a quantity to be traded against other quantities \parencite{kant1998groundwork}. An optimizing frame can represent such a requirement only by falsifying it: either as a very large term that a still larger term could in principle outweigh, which denies its unconditionality, or as an absolute constraint bounding the feasible set, which removes it from the objective altogether and concedes that the moral was never a matter of maximization to begin with. The second feature is the \emph{non-tradability} of moral considerations. To make a duty a term in an objective is to make it tradable, available to be sacrificed when the numbers favor the trade; but much of what we recognize as moral seriousness consists precisely in the refusal to make certain considerations tradable, in holding that there are things one does not do whatever the yield, and this refusal is unintelligible within a frame in which everything is, by the form of the objective, a quantity that trades \parencite{williams1985limits}. The man in the case did the tradable thing; that the friend's dependence could be traded away for a sufficient gain is exactly what a moral seriousness would deny. The third feature is the \emph{irreplaceability of the other}. The person to whom one is bound is not, morally, the bearer of a quantity of value that another bearer of equal quantity could replace without loss; the other is irreplaceable, and the claim the other makes is the claim of this singular one and not of the value it happens to carry \parencite{levinas1969totality}. Optimization, whose feasible set is a space of interchangeable options ranked by a common measure, cannot represent irreplaceability at all, since its every option is in principle substitutable by another of equal or greater value; and irreplaceability is precisely the denial that such substitution is ever, here, without loss.

\begin{claim}[Orthogonality of the optimal and the good]
The criterion of the \emph{optimal} and the criterion of the \emph{good} are orthogonal: the first asks whether an option maximizes a given objective, the second whether it ought to be chosen, and success on the first is no success on the second. An optimizing practice, even one well-posed and reaching a genuine maximum, may therefore fail as an ethical practice, and does so wherever the good it maximizes is not the good it ought to honor. Morality resists the optimizing form in particular because of the unconditionality of its requirements, the non-tradability of its considerations, and the irreplaceability of the other, none of which an objective function can represent without falsifying. Morality is therefore not a higher-order optimization, a maximization of some moral quantity, but a practice outside the grammar of optimization. The optimal is not the good.
\end{claim}

\subsection*{The completion of the criticism}

With this the criticism completes itself, and it does so by closing the last exit the optimizing frame might take. Confronted with the argument thus far, a defender of optimization has one sophisticated move left: to concede that ordinary utility is the wrong objective and to propose that morality itself be optimized, that we maximize not pleasure or preference but the good, folding the moral considerations into a richer objective and maximizing that. This section's claim is what forecloses the move. Morality cannot be the objective of an optimization, because its central features, the unconditional, the non-tradable, the irreplaceable, are exactly the features an objective function cannot carry; to write morality into an objective is to strip it of these features and so to write in something else under its name. The optimizing frame cannot be rescued by aiming it at the good, for the good is not the kind of thing there is a maximum of. And this does not leave ethics without rationality, which would be a disaster of the opposite kind and which the paper is at pains to avoid. It leaves ethics with a rationality that is not optimization, a way of reasoning about what to do that is answerable, discriminating, and severe without being the maximization of any measure; and the recovery of that other rationality, promised from the start, is the work of the part that follows.
